\section{Introduction}
\sloppy

Forecasting how traffic or rainfall evolves over a physical network is a long-standing data-mining task with two
competing recipes. Specialists, mostly graph neural networks and transformers, are trained on one network's
history~\citep{li2018dcrnn,wu2019graphwavenet,shao2022stid,liu2023staeformer}. Time-series foundation models such
as TimesFM~\citep{das2024timesfm} are pretrained once on large corpora and forecast any series zero-shot. Each
recipe is usually evaluated in its own protocol: specialists on point error over fixed windows, often with numbers
copied from earlier papers, and foundation models on aggregate leaderboards of unrelated series. A practitioner
choosing between them for a given network has no like-for-like evidence, and benchmark studies have shown that
re-running published specialists under one pipeline changes their rankings~\citep{shao2024basicts}.

We ask a narrow question with a fixed protocol: on real sensor networks, how does a small model trained from
scratch on the network's own history compare with (a)~published specialists run with their authors' code and
(b)~zero-shot foundation models, on the same windows, with the same point and probabilistic scores, when a third of
the sensors go dark, and with tests that say whether the differences are larger than noise?

\paragraph{Contributions.}
\begin{itemize}
\item \textbf{PWM-Dyn} (\S\ref{sec:model}), a $3\times10^5$-parameter graph-diffusion model with explicit missing-data
inputs, three quantile heads and split-conformal calibration, trained in 9--52 GPU-minutes per seed on a free T4.
\item \textbf{\benchname{}} (\S\ref{sec:bench}), an open benchmark of six traffic networks and two monsoon rainfall
regions with four tracks (standard context, open context, dead sensors, zero-shot transfer), one scorer for point,
quantile and interval quality, saved forecasts and paired significance tests.
\item \textbf{A controlled comparison} (\S\ref{sec:exp}) against STID, Graph WaveNet and STAEformer trained with
their official code and published hyper-parameters on our windows, and against TimesFM 2.5 and 3.0 zero-shot. Every
PWM-Dyn configuration choice was made on validation data and the reported models use seeds never used during
development; four of the six traffic networks were never touched during development.
\item \textbf{Findings, including losses.} Under a 12-step context a three-seed PWM-Dyn ensemble has the lowest MAE of the four trained models on five of six traffic networks and is
second on PEMS08, where STAEformer leads by 1.4\% even though its run was cut short; against STID, the
closest rival, the margins are significant on three networks and not significant on PEMS08, PEMS04
and PEMS07,
and single PWM-Dyn seeds are slightly worse than single STID runs on three networks. Against zero-shot
TimesFM~3.0 it has lower MAE on seven of eight tasks (rainfall entries from the previous version) and lower interval score on seven of eight, and loses on PEMS03.
With 30\% of sensors dark it beats per-sensor climatology on seven of eight tasks while specialists without a
missing-data mechanism fail, but climatology beats it on PEMS07. An ablation shows that a climatology input added in an earlier
version cost accuracy under the short context; removing it, chosen on validation, gave the reported model.

\end{itemize}
